% Wersja źródłowa: audyt IV, jednostka 07
\setcounter{section}{6}
\section{Oddzielanie zbiorów wypukłych, polary i twierdzenie bipolarne}\label{sec:separacja-bipolarne}

\subsection{Przejście do rzeczywistych funkcjonałów}

W \hyperref[sec:hahn-banach]{rozdziale 6} udowodniliśmy twierdzenie Hahna--Banacha dla funkcjonałów liniowych.
Teraz wykorzystamy je do opisu geometrii zbiorów wypukłych i topologii,
którą można odtworzyć za pomocą przestrzeni sprzężonej. W szczególności odróżnimy
\emph{nierówności ścisłe dla każdego punktu} od \emph{oddzielenia z jednolitą luką}.

Przez cały rozdział $\K\in\{\R,\C\}$, a $X$ jest przestrzenią liniowo-topologiczną
nad $\K$. Gdy jest to wyraźnie zaznaczone, zakładamy lokalną wypukłość
oraz własność Hausdorffa. W topologii zespolonej oddzielanie zbiorów
odbywa się za pomocą \emph{rzeczywistych części} funkcjonałów z $X'$.

\begin{definicja}[Rzeczywiste funkcjonały i rodzaje separacji]\label{def:separacja}
Niech $A,B\subseteq X$ będą niepustymi zbiorami, a
$\ell:X\to\R$ niezerowym, ciągłym, rzeczywiście liniowym funkcjonałem
(na przestrzeni $X$ rozpatrywanej nad $\R$). Mówimy, że
\emph{oddziela} $A$ i $B$, gdy
$\sup_{a\in A}\ell(a)\le\inf_{b\in B}\ell(b)$.
Oddziela je \emph{ściśle punktowo}, gdy $\ell(a)<\ell(b)$
dla wszystkich $a\in A$, $b\in B$, a \emph{silnie}, gdy istnieje
$\delta>0$ takie, że $\sup_A\ell+\delta\le\inf_B\ell$.
Hiperpłaszczyzną afiniczną funkcjonału $\ell$ na poziomie $c\in\R$
nazywamy $H_{\ell,c}=\{x\in X:\ell(x)=c\}$.
\end{definicja}

\begin{lemat}[Rzeczywista część funkcjonału zespolonego]\label{lem:realification}
Niech $X$ będzie zespoloną TVS. Każdy ciągły, rzeczywiście liniowy
$\ell:X\to\R$ ma dokładnie jedno przedstawienie
$\ell=\operatorname{Re}F$, gdzie $F\in X'$. Mianowicie
$F(x)=\ell(x)-i\ell(ix)$.
\end{lemat}
\begin{proof}
Liniowość zespolona i jednoznaczność zostały wykazane w
lemacie o odtworzeniu funkcjonału zespolonego z \hyperref[sec:hahn-banach]{rozdziału 6}. Ponieważ $x\mapsto ix$
jest ciągły, także funkcjonał $F$ jest ciągły.
\end{proof}

\subsection{Oddzielanie zbioru otwartego od punktu}

\begin{lemat}[Funkcjonał Minkowskiego dla zbioru wypukłego]\label{lem:open-gauge}
Niech $E$ będzie rzeczywistą TVS, a $V\subseteq E$ otwartym
wypukłym otoczeniem zera (niekoniecznie symetrycznym).
Wówczas $V$ jest pochłaniający, a jego funkcjonał
$p_V(x)=\inf\{t>0:x\in tV\}$ jest skończony, podliniowy i ciągły.
Ponadto $V=\{x:p_V(x)<1\}$.
\end{lemat}
\begin{proof}
Pochłanianie wynika z ciągłości $t\mapsto tx$ w zerze.
Dla dowolnego $x\in E$ i $r>p_V(x)$ zachodzi $x\in rV$:
istnieje bowiem $0<s<r$ z $x\in sV$, a ponieważ $0\in V$
i $V$ jest wypukły, mamy $(s/r)V\subseteq V$.
Dla $t>0$ bezpośrednio z definicji otrzymujemy
$p_V(tx)=t p_V(x)$, a $p_V(0)=0$. Jeżeli
$r>p_V(x)$ i $s>p_V(y)$, to możemy napisać
$x=ru$, $y=sv$ z $u,v\in V$. Wypukłość daje
$(x+y)/(r+s)=(r/(r+s))u+(s/(r+s))v\in V$,
czyli $p_V(x+y)\le r+s$. Przechodząc z $r$ i $s$
do odpowiednich infimów, otrzymujemy
$p_V(x+y)\le p_V(x)+p_V(y)$. Zatem $p_V$ jest
podliniowy; w tym dowodzie nie korzystaliśmy z symetrii $V$.
Z wypukłości i $0\in V$ wynika
$tV\subseteq V$ dla $0<t\le1$. Jeżeli $x\in V$, to
z otwartości $V$ i ciągłości $s\mapsto sx$ w $s=1$
można wybrać $s>1$ takie, że $sx\in V$, czyli
$p_V(x)\le1/s<1$. Odwrotnie, jeżeli $p_V(x)<1$,
można znaleźć $t<1$ z $x\in tV\subseteq V$.
Stąd $V=\{p_V<1\}$.

Dla $\varepsilon>0$ zbiór $\{p_V<\varepsilon\}=\varepsilon V$
jest otwarty. Zatem $p_V$ jest górnie półciągły. Aby
wykazać ciągłość, dla każdego $x,y\in E$ korzystamy z
$\,p_V(x)\le p_V(y)+p_V(x-y)$ i
$p_V(y)\le p_V(x)+p_V(y-x)$.
Jeżeli $x-y\in\varepsilon V\cap(-\varepsilon V)$,
to obie różnice są mniejsze niż $\varepsilon$.
Takie zbiory tworzą otoczenia zera, co daje ciągłość.
\end{proof}

\begin{twierdzenie}[Oddzielanie punktu od otwartego zbioru wypukłego]
\label{thm:open-point-sep}
Niech $X$ będzie lokalnie wypukłą TVS nad $\K$,
$C\subseteq X$ niepustym otwartym zbiorem wypukłym
(w sensie rzeczywistym), a $x_0\in X\setminus C$.
Wówczas istnieje niezerowy ciągły, rzeczywiście liniowy
$\ell:X\to\R$ taki, że $\ell(c)<\ell(x_0)$
dla każdego $c\in C$. Jeżeli $\K=\C$, można przyjąć
$\ell=\operatorname{Re}F$ z $F\in X'$.
Nie twierdzimy ogólnie, że $\sup_C\ell<\ell(x_0)$.
\end{twierdzenie}
\begin{proof}
Traktujemy $X$ jako rzeczywistą TVS; lokalna wypukłość
jest zachowana. Zbiór $D=C-x_0$ jest niepusty, otwarty
i wypukły, a $0\notin D$. Wybierzmy $a\in D$ i połóżmy
$V=D-a$. Zbiór $V$ jest otwartym, wypukłym otoczeniem zera,
a $-a\notin V$, stąd $p_V(-a)\ge1$ na mocy
Lematu~\ref{lem:open-gauge}. Na prostej $\R(-a)$
określmy $g(t(-a))=t$. Dla $t\ge0$ mamy
$g(t(-a))=t\le p_V(t(-a))$. Dla $t<0$ nierówność
także zachodzi, bo $p_V\ge0$. Rzeczywiste twierdzenie
Hahna--Banacha daje liniowe $\ell:X\to\R$ spełniające
$\ell\le p_V$ i $\ell(-a)=1$. Z ciągłości $p_V$
i nierówności $-p_V(-x)\le\ell(x)\le p_V(x)$
wynika ciągłość $\ell$.

Dla dowolnego $d\in D$ mamy $d-a\in V$, więc
$\ell(d-a)\le p_V(d-a)<1$. Ponieważ $\ell(a)=-1$,
otrzymujemy $\ell(d)<0$, czyli
$\ell(c)<\ell(x_0)$ dla $c\in C$. Równość
$\ell(-a)=1$ dowodzi, że $\ell\ne0$. W przypadku
zespolonym stosujemy Lemat~\ref{lem:realification}.
\end{proof}

\begin{twierdzenie}[Oddzielanie dwóch zbiorów wypukłych]
\label{thm:open-two-sep}
Niech $X$ będzie lokalnie wypukłą TVS nad $\K$,
a $A,B\subseteq X$ niepustymi, rozłącznymi zbiorami wypukłymi.
Jeśli $A$ jest otwarty, istnieje niezerowy, ciągły, rzeczywiście
liniowy funkcjonał $\ell$ spełniający
$\ell(a)<\ell(b)$ dla wszystkich $a\in A$, $b\in B$.
W szczególności istnieje $c\in\R$ takie, że
$\sup_A\ell\le c\le\inf_B\ell$.
\end{twierdzenie}
\begin{proof}
Zbiór $D=A-B=\{a-b:a\in A,\ b\in B\}$ jest wypukły
i otwarty (jest sumą zbioru otwartego i dowolnego),
a $0\notin D$. Z Twierdzenia~\ref{thm:open-point-sep}
dostajemy $\ell(d)<0$ dla wszystkich $d\in D$,
co daje nierówność punktową. Dla wybranego $b_0\in B$
mamy $\sup_A\ell\le\ell(b_0)<\infty$; analogicznie,
dla $a_0\in A$ mamy $\inf_B\ell\ge\ell(a_0)>-\infty$.
Stąd obie liczby są skończone i możemy wybrać $c$
pomiędzy nimi.
\end{proof}

\begin{przyklad}[Dlaczego otwartość nie daje jednolitej luki?]
Dla $X=\R$, $A=(-1,0)$, $B=\{0\}$ funkcjonał
$\ell(x)=x$ spełnia $\ell(a)<\ell(b)$ dla wszystkich
$a\in A$, $b\in B$, lecz $\sup_A\ell=\inf_B\ell=0$.
Nie istnieje silna separacja tych zbiorów, ponieważ
$0\in\overline A\cap B$.
\end{przyklad}

\subsection{Oddzielanie silne i hiperpłaszczyzny podpierające}

\begin{lemat}[Wspólne otoczenie dla zbioru zwartego]
\label{lem:compact-neighborhood}
Niech $X$ będzie TVS, $K\subseteq X$ zbiorem zwartym,
a $F\subseteq X$ zbiorem domkniętym, przy czym $K\cap F=\varnothing$.
Wtedy istnieje otoczenie $U$ zera takie, że
$(K+U)\cap F=\varnothing$. Jeżeli $X$ jest lokalnie wypukła,
można przyjąć $U$ otwarte i wypukłe.
\end{lemat}
\begin{proof}
Dla każdego $k\in K$ istnieje otoczenie $V_k$ zera
z $(k+V_k)\cap F=\varnothing$. Wybieramy otoczenie
$W_k$ zera takie, że $W_k+W_k\subseteq V_k$,
a w przypadku lokalnie wypukłym możemy żądać
jego wypukłości i otwartości. Zwartość $K$ daje
pokrycie skończone $K\subseteq\bigcup_{j=1}^m(k_j+W_{k_j})$.
Wybierzmy otoczenie $U\subseteq\bigcap_j W_{k_j}$.
Jeżeli $x\in K$ należy do $k_j+W_{k_j}$, to
$x+U\subseteq k_j+V_{k_j}$, zatem $(K+U)\cap F=\varnothing$.
\end{proof}

\begin{twierdzenie}[Silna separacja zbioru zwartego i domkniętego]
\label{thm:strong-separation}
Niech $X$ będzie lokalnie wypukłą TVS nad $\K$,
$A,B\subseteq X$ niepustymi, rozłącznymi zbiorami wypukłymi.
Jeżeli $A$ jest zwarty, a $B$ domknięty, to istnieje
niezerowy ciągły, rzeczywiście liniowy funkcjonał $\ell$
oraz $\delta>0$ takie, że
\[
\sup_{a\in A}\ell(a)+\delta\le\inf_{b\in B}\ell(b).
\]
W przypadku zespolonym $\ell=\operatorname{Re}F$ dla pewnego $F\in X'$.
\end{twierdzenie}
\begin{proof}
Na mocy Lematu~\ref{lem:compact-neighborhood} istnieje
otwarte wypukłe otoczenie $U$ zera takie, że
$(A+U)\cap B=\varnothing$. Zbiór $A+U$ jest otwarty
i wypukły. Z Twierdzenia~\ref{thm:open-two-sep}
otrzymujemy niezerowy $\ell$ z
$\ell(a+u)<\ell(b)$ dla wszystkich
$a\in A$, $u\in U$, $b\in B$. Ponieważ $\ell\ne0$
i $U$ jest otoczeniem zera, istnieje $u_0\in U$
z $\ell(u_0)>0$. Kładąc $\delta=\ell(u_0)$, uzyskujemy
$\ell(a)+\delta<\ell(b)$ dla wszystkich $a,b$,
a zatem $\sup_A\ell+\delta\le\inf_B\ell$.
\end{proof}

\begin{wniosek}[Separacja punktu i domkniętego zbioru wypukłego]
\label{cor:point-closed-convex}
Niech $X$ będzie lokalnie wypukłą TVS, $C\subseteq X$
niepustym, domkniętym i wypukłym zbiorem oraz
$x_0\in X\setminus C$. Wówczas istnieją
ciągły rzeczywiście liniowy $\ell\ne0$ i $\delta>0$
takie, że $\ell(x_0)+\delta\le\inf_{c\in C}\ell(c)$.
\end{wniosek}
\begin{proof}
Stosujemy Twierdzenie~\ref{thm:strong-separation}
do $A=\{x_0\}$ i $B=C$.
\end{proof}

\begin{twierdzenie}[Hiperpłaszczyzna podpierająca]
\label{thm:supporting-hyperplane}
Niech $X$ będzie lokalnie wypukłą TVS nad $\K$,
a $C\subseteq X$ niepustym, domkniętym i wypukłym
zbiorem o niepustym wnętrzu. Dla każdego $x_0\in\partial C$
istnieje niezerowy, ciągły, rzeczywiście liniowy
$\ell:X\to\R$ taki, że $\ell(x)\le\ell(x_0)$
dla wszystkich $x\in C$.
\end{twierdzenie}
\begin{proof}
Wnętrze $G=\operatorname{int} C$ jest wypukłe.
Ponieważ $x_0\in\partial C$, $x_0\notin G$.
Twierdzenie~\ref{thm:open-point-sep} daje $\ell\ne0$
z $\ell(g)<\ell(x_0)$ dla $g\in G$.
Dla $x\in C$ i ustalonego $g_0\in G$ punkty
$x_t=(1-t)x+tg_0$ należą do $G$ dla $0<t\le1$:
małe otoczenie zera $V$ z $g_0+V\subseteq C$
spełnia $x_t+tV\subseteq C$ na mocy wypukłości.
Ponieważ $x_t\to x$ przy $t\downarrow0$, ciągłość
$\ell$ daje $\ell(x)\le\ell(x_0)$.
\end{proof}

\begin{uwaga}[Czego nie należy wnioskować]
Dla zbioru domkniętego wypukłego o pustym wnętrzu
funkcjonał podpierający w zadanym punkcie nie musi
istnieć bez dodatkowych hipotez. Podobnie z istnienia
hiperpłaszczyzny podpierającej nie wynika, że jest ona
jednoznaczna. W przestrzeni zespolonej pojęcie
hiperpłaszczyzny podpierającej odnosi się do równania
$\operatorname{Re}F(x)=c$, a nie do $F(x)=c$.
\end{uwaga}

\subsection{Zbiory absolutnie wypukłe i polary}

\begin{definicja}[Otoczka absolutnie wypukła]
\label{def:absolutely-convex-hull}
Niech $X$ będzie przestrzenią liniową nad $\K$, a $A\subseteq X$.
Definiujemy
\[
\operatorname{aco}(A)=\left\{\sum_{j=1}^n\lambda_j a_j:
 n\in\N,\ a_j\in A,\ \lambda_j\in\K,
 \sum_{j=1}^n|\lambda_j|\le1\right\}\cup\{0\}.
\]
Jest to najmniejszy zbiór wypukły i zrównoważony zawierający $A$.
Przyjmujemy $\operatorname{aco}(\varnothing)=\{0\}$.
\end{definicja}

\begin{definicja}[Polary w parze $(X,X')$]\label{def:polar}
Niech $X$ będzie TVS nad $\K$ i $X'$ przestrzenią sprzężoną.
Dla $A\subseteq X$ i $B\subseteq X'$ definiujemy
\[
 \begin{aligned}
 A^\circ&=\{f\in X':|f(a)|\le1\ \text{dla każdego }a\in A\},\\
 B^\circ&=\{x\in X:|f(x)|\le1\ \text{dla każdego }f\in B\}.
 \end{aligned}
\]
W szczególności $A^{\circ\circ}=(A^\circ)^\circ\subseteq X$.
Powyższa biegunowość jest \emph{absolutna}: w przypadku
zespolonym również używamy modułu, nie samej części rzeczywistej.
\end{definicja}

\begin{stwierdzenie}[Podstawowe własności polary]\label{prop:polar-properties}
Niech $X$ będzie TVS nad $\K$ i $A,A_1,A_2\subseteq X$.
Wówczas:
\begin{enumerate}
\item $A_1\subseteq A_2$ implikuje $A_2^\circ\subseteq A_1^\circ$;
\item $A^\circ=(\operatorname{aco}A)^\circ$ oraz
      $A\subseteq A^{\circ\circ}$;
\item $A^\circ$ jest wypukły, zrównoważony i domknięty
      w topologii punktowej zbieżności na $X'$, tj. $\sigma(X',X)$;
\item $A^{\circ\circ}$ jest wypukły, zrównoważony i domknięty
      w oryginalnej topologii $X$.
\end{enumerate}
\end{stwierdzenie}
\begin{proof}
Pierwsza własność wynika bezpośrednio z kwantyfikatora
w definicji. Dla $f\in A^\circ$ i
$x=\sum_j\lambda_j a_j$ z $\sum_j|\lambda_j|\le1$
mamy $|f(x)|\le\sum_j|\lambda_j||f(a_j)|\le1$,
co dowodzi równości polar. Każdy $a\in A$ spełnia
$|f(a)|\le1$ dla $f\in A^\circ$, zatem $a\in A^{\circ\circ}$.
Zbiory $\{f:|f(a)|\le1\}$ są domknięte, wypukłe i
zrównoważone w topologii punktowej; ich przecięcie
jest zbiorem $A^\circ$. Analogicznie przecięcie
$\{x:|f(x)|\le1\}$ po $f\in A^\circ$ daje własność (4).
\end{proof}

\begin{twierdzenie}[Twierdzenie bipolarne]\label{thm:bipolar}
Niech $X$ będzie lokalnie wypukłą TVS Hausdorffa nad
$\K\in\{\R,\C\}$, a $A\subseteq X$ dowolnym zbiorem.
Wówczas
\[
 A^{\circ\circ}=\overline{\operatorname{aco}(A)},
\]
gdzie domknięcie jest brane w oryginalnej topologii $X$.
\end{twierdzenie}
\begin{proof}
Z poprzedniego stwierdzenia wynika
$\overline{\operatorname{aco}(A)}\subseteq A^{\circ\circ}$.
Oznaczmy $C=\overline{\operatorname{aco}(A)}$.
Jest to zbiór niepusty, domknięty, wypukły i zrównoważony.
Jeśli $x_0\notin C$, Wniosek~\ref{cor:point-closed-convex}
zastosowany do $C$ daje ciągły rzeczywiście liniowy
$\ell$ z $\ell(x_0)>\sup_{c\in C}\ell(c)$
(po ewentualnej zmianie znaku).
Niech $F=\ell$ w przypadku rzeczywistym, a
$F(x)=\ell(x)-i\ell(ix)$ w zespolonym.
Wówczas $F\in X'$, $\operatorname{Re}F=\ell$.
Ze zrównoważenia $C$ wynika
\[
 m:=\sup_{c\in C}|F(c)|=\sup_{c\in C}\operatorname{Re}F(c)
 =\sup_{c\in C}\ell(c)<\ell(x_0)\le|F(x_0)|.
\]
W szczególności $m$ jest skończone i nieujemne.
Jeżeli $m>0$, to $G=F/m$ spełnia $|G(a)|\le1$
dla wszystkich $a\in A$, lecz $|G(x_0)|>1$.
Jeżeli $m=0$, wybieramy dodatnie $t$ z
$t|F(x_0)|>1$ i kładziemy $G=tF$.
W obu przypadkach $G\in A^\circ$, a $x_0\notin A^{\circ\circ}$.
Otrzymujemy więc odwrotną inkluzję.
\end{proof}

\begin{wniosek}[Domknięcie wypukłe nie zmienia się po osłabieniu topologii]
\label{cor:weak-convex-closure}
Niech $X$ będzie lokalnie wypukłą TVS Hausdorffa,
$C\subseteq X$ zbiorem wypukłym, a $\sigma(X,X')$
najsłabszą topologią, w której wszystkie $f\in X'$
są ciągłe. Wtedy domknięcie $C$ w $\sigma(X,X')$
jest równe domknięciu $C$ w oryginalnej topologii $X$.
\end{wniosek}
\begin{proof}
Jeśli $C=\varnothing$, teza jest oczywista. W dalszej części
zakładamy $C\ne\varnothing$. Topologia $\sigma(X,X')$ jest słabsza od wyjściowej,
zatem domknięcie słabe zawiera domknięcie wyjściowe.
Niech $D=\overline C$ oznacza domknięcie w topologii
wyjściowej. Zbiór $D$ jest niepusty, domknięty i wypukły.
Jeżeli $x_0\notin D$, stosujemy
Wniosek~\ref{cor:point-closed-convex} do $D$, a nie do
niekoniecznie domkniętego $C$. Otrzymujemy ciągły
rzeczywiście liniowy $\ell$ i $\alpha\in\R$ takie, że
$\ell(x_0)<\alpha<\inf_{d\in D}\ell(d)
\le\inf_{c\in C}\ell(c)$.
Jego reprezentacja przez część rzeczywistą funkcjonału
$F\in X'$ pokazuje, że zbiór
$\{x:\operatorname{Re}F(x)<\alpha\}$ jest słabo otwartym
otoczeniem $x_0$ rozłącznym z $C$.
\end{proof}

\subsection{Przykłady i pułapki}

\begin{przyklad}[Polary kul jednostkowych w $\R^2$]\label{ex:polar-balls}
W parze $X=X'=\R^2$ z parowaniem
$\langle a,x\rangle=a_1x_1+a_2x_2$ polara
$A=\{x:|x_1|+|x_2|\le1\}$ jest zbiorem
$A^\circ=\{a:\max(|a_1|,|a_2|)\le1\}$.
Jest to konsekwencja tożsamości
$\sup_{\|x\|_1\le1}|\langle a,x\rangle|=\|a\|_\infty$.
\begin{center}
\begin{tikzpicture}[scale=1.07,>=Stealth]
\begin{scope}[shift={(-2.65,0)}]
\draw[->,gray] (-1.55,0)--(1.55,0) node[right] {$x_1$};
\draw[->,gray] (0,-1.48)--(0,1.5) node[above] {$x_2$};
\fill[blue!12] (1,0)--(0,1)--(-1,0)--(0,-1)--cycle;
\draw[blue!65!black,very thick] (1,0)--(0,1)--(-1,0)--(0,-1)--cycle;
\node at (0,-1.85) {$A:\ \|x\|_1\le1$};
\end{scope}
\draw[->,thick] (-.72,0)--(.75,0) node[midway,above] {$\circ$};
\begin{scope}[shift={(2.65,0)}]
\draw[->,gray] (-1.55,0)--(1.55,0) node[right] {$a_1$};
\draw[->,gray] (0,-1.48)--(0,1.5) node[above] {$a_2$};
\fill[blue!12] (-1,-1) rectangle (1,1);
\draw[blue!65!black,very thick] (-1,-1) rectangle (1,1);
\node at (0,-1.85) {$A^\circ:\ \|a\|_\infty\le1$};
\end{scope}
\end{tikzpicture}
\end{center}
\end{przyklad}

\begin{przyklad}[Polara w produkcie $\K^\N$]\label{ex:polar-product}
Niech $X=\K^\N$ z topologią produktową, a
$A=\{e_n:n\in\N\}$, gdzie $e_n$ ma jedynkę na
$n$-tej współrzędnej. Z \hyperref[sec:hahn-banach]{rozdziału 6} wiemy, że $X'$
składa się z funkcjonałów $f_a(x)=\sum_n a_nx_n$
o skończonym nośniku. Zatem
$A^\circ=\{f_a:\ a\in c_{00},\ |a_n|\le1\ \text{dla każdego }n\}$,
a więc
\[
 A^{\circ\circ}=
 \left\{x\in\K^\N:\sum_{n=1}^\infty|x_n|\le1\right\}.
\]
Istotnie, supremum $|\sum_n a_nx_n|$ po
skończenie wspartych ciągach $(a_n)$ z $|a_n|\le1$
wynosi $\sup_{m\in\N}\sum_{n=1}^m|x_n|$. Z drugiej strony $\operatorname{aco}(A)$ składa się dokładnie
z ciągów o skończonym nośniku spełniających
$\sum_n|x_n|\le1$. Dla każdego $m$ funkcja
$x\mapsto\sum_{n=1}^m|x_n|$ jest ciągła w topologii
produktowej i nie przekracza $1$ na $\operatorname{aco}(A)$.
Dlatego nie przekracza też $1$ na jego domknięciu,
a po wzięciu supremum po $m$ dostajemy
$\sum_{n=1}^\infty|x_n|\le1$. Odwrotnie, jeśli
$x$ spełnia tę nierówność, jego obcięcia
$x^{(m)}=(x_1,\ldots,x_m,0,\ldots)$ należą do
$\operatorname{aco}(A)$ i zbiegają do $x$ w topologii
produktowej. Dowodzi to bezpośrednio opisanego
w twierdzeniu bipolarnym domknięcia.
\end{przyklad}

\begin{przyklad}[Dlaczego lokalna wypukłość ma znaczenie?]
Niech $X=L^p([0,1])$, $0<p<1$, z naturalną topologią
metryczną. Jak wiemy, $X'=\{0\}$. Dlatego dla
$A=\{0\}$ zachodzi $A^\circ=\{0\}$ oraz
$A^{\circ\circ}=X$, mimo że
$\overline{\operatorname{aco}(A)}=\{0\}$.
Twierdzenie bipolarne nie zachodzi więc w dowolnych TVS.
\end{przyklad}

\begin{uwaga}[Co dalej?]
W następnym rozdziale szczegółowo zbadamy pary
przestrzeni sprzężonych i topologie $\sigma(X,X')$ oraz
$\sigma(X',X)$. Dzięki temu pojęcia polary i
zamknięcia słabego otrzymają naturalną interpretację
jako narzędzia opisu przestrzeni funkcjonałów ciągłych,
a później także dystrybucji.
\end{uwaga}
